%% kapa: exam=1 round=21 year=2010 period=1 form=A subject=civil_law subject_ko=민법 seq=18 number=18
%% source: Q-Net 감정평가사 기출문제 (cst00309) artlSeq=5024181 file=extracted/1008/5024181/2010년 제21회 감정평가사 1차시험 기출문제/1교시.hwp
%% conversion: kordoc 4.24.51 -> tools/gen1.py (auto); review=proofread-sonnet; flags=none
\begin{question}{18}
선의취득에 관한 설명으로 옳지 않은 것은? (다툼이 있으면 판례에 의함)
\begin{choices}
  \item 선의취득에서 선의⋅무과실의 기준시점은 물권행위가 완성되는 시점이므로, 인도가 물권적 합의보다 먼저 이루어진 경우에는 물권적 합의시를 기준으로 판단하여야 한다.
  \item 상속으로 동산의 점유를 취득한 경우에는 선의취득이 인정될 수 없다.
  \item 채무자 이외의 자의 소유에 속하는 동산을 경매절차에서 매각대금을 납부하고 인도받은 매수인은 그 동산을 선의취득할 수 없다.
  \item 저당권이나 지상권과 같이 부동산에 관한 권리는 선의취득의 객체가 되지 못한다.
  \item 점유개정의 방법으로 동산에 대한 이중의 양도담보설정계약이 체결된 경우, 뒤에 계약을 체결한 후순위 채권자는 양도담보권을 선의취득할 수 없다.
\end{choices}
\end{question}
\answer{3}
